\section{Conclusion}
\label{sec:conclusion}

In this work, we explored the concept of origin 'mode' as a unified latent variable for text source characterization.
Through experiments with \aigen, \dictated, and \typed text, we demonstrated that while 'mode' is inherently captured in text embeddings with moderate separation (Silhouette Score: \silhouette), current LLMs struggle to zero-shot distinguish between fine-grained human artifacts like phonetic dictation errors and orthographic typing mistakes.
The high recall for \aigen text confirms that AI-generated signatures are highly detectable, but the confusion between human modes suggests that artifact granularity remains a significant challenge.

{\bf Key Takeaway.}
Our findings validate that text origin 'mode' acts as a structured latent variable, but its explicit isolation in zero-shot contexts requires models to overcome their internal language priors that tend to 'gloss over' systematic human errors.

{\bf Future Work.}
Future research should focus on leveraging organic ASR and typing datasets to capture more complex real-world artifacts.
Additionally, we propose investigating supervised disentanglement techniques, such as FactorVAE, to explicitly separate the 'mode' dimension from content and style in embedding spaces.
Extending this analysis to few-shot prompting and larger-scale models will also be crucial for reaching human-level granularity in origin mode isolation.
